\Needspace{8\baselineskip}\originalchapter{2006 年期末卷}
题源: 19 Dec 2006, 9--12时, 原卷8页、4题各25\%. 当期官方解答未取得. 原扫描公式已人工复核, 纠正文本抽取的系数倒序.
\begin{exercise}\label{ex:06-final-1}
原题1. 轴对称环形体系, 导引中心极向坐标x=R\ensuremath{-}R0,y=Z或r,\ensuremath{\theta}. (a)纯环场 $B_\varphi=B_0R_0/R$, 初$x=y=0$, 动能W0=W\ensuremath{\perp}0+W\ensuremath{\parallel}0, 求x(t),y(t). (b)加入弱 $B_\theta=B_{\theta0}r/R$, 忽略漂移, 初$r=r_0,\theta=0$, 求能绕极向2\ensuremath{\pi}的最大v\ensuremath{\perp}0/v0. (c)纯平行v\ensuremath{\perp}0=0, $\dot y=\omega x,\dot x=-\omega y$, 求\ensuremath{\omega}和轨道; (d)加纯环场估计的漂移, 改方程; (e) $B_0^2\gg B_{\theta0}^2$使\ensuremath{\omega}约常数, 初$r=r_0,\theta=0$, 求f(x,y)=0. 原(e)引用(c), 按(d)有漂移轨道解, 同时给(c)无漂移极限.
\end{exercise}
\begin{solution}
(a)曲率/gradB均沿+Z(对q正、B\ensuremath{\phi}正), $v_D=(W_\perp+2W_\parallel)/(qBR)=(W_{\perp0}+2W_{\parallel0})/(qB_0R_0)$, 因初$r=0$且无径向漂移, $x=0,y=v_Dt$. (b)最弱场外侧R0+r0, 最强内侧R0\ensuremath{-}r0, 可通过要求 $v_{\perp0}/v_0<\sqrt{(R_0-r_0)/(R_0+r_0)}$, 等号为临界镜点, 不保证有限时间完成2\ensuremath{\pi}. (c) $\omega=v_\parallel B_{\theta0}/(B_0R_0)$, x=r0cos\ensuremath{\omega}t,y=r0sin\ensuremath{\omega}t, 圆轨道.
(d)纯平行 $v_D=mv_\parallel^2/(qB_0R_0)$, $\dot x=-\omega y$, $\dot y=\omega x+v_D$. (e)令d=vD/\ensuremath{\omega}, 解 $x+d=(r_0+d)\cos\omega t$, $y=(r_0+d)\sin\omega t$, 所以 $(x+d)^2+y^2=(r_0+d)^2$. 电荷改变d方向, 轨道中心偏移; vd0退化(c). 纯平行且弱极向场的近似不包含有限\ensuremath{\mu}粒子的banana轨道.
\begin{center}\begin{tikzpicture}[scale=.75]\draw[->](-3,0)--(3,0)node[right]{$x$};\draw[->](0,-2.5)--(0,2.5)node[above]{$y$};\draw[dashed](0,0)circle(1.4);\draw[thick,structurecolor](-.5,0)circle(1.9);\fill(-.5,0)circle(2pt)node[below]{$-d$};\fill(1.4,0)circle(2pt)node[above]{$r_0$};\node at(0,-3){同一初点, 虚线无漂移、实线有漂移};\end{tikzpicture}\end{center}
\end{solution}
\begin{exercise}\label{ex:06-final-2}
原题2. 由给表估热扩散, 表n用m$^{-3}$,T用keV,B用T, ln\ensuremath{\Lambda}15. 电子 $\nu_e=1.38\times10^{-15}n/T_e^{3/2}$, $v_{te}=1.33\times10^7\sqrt{T_e}$m/s, $\rho_e=1.07\times10^{-4}\sqrt{T_e}/B$m; 质子 $\nu_i=2.5\times10^{-17}n/T_i^{3/2}$, $v_{ti}=3.09\times10^5\sqrt{T_i}$m/s, $\rho_i=4.57\times10^{-3}\sqrt{T_i}/B$m; \ensuremath{\eta}=$2.5\times10^{-8}T_e^{-3/2}$\ensuremath{\Omega}m. (a)横\ensuremath{\chi}; (b)平行\ensuremath{\chi}; (c)宽2\ensuremath{\ell}的无限平板、B/J沿z,n/B/J均匀,两物种T相同, $-\nabla\cdot q+\eta J^2=0$, $q_y=-\alpha T^{-1/2}T'$、\ensuremath{\alpha}常数, \ensuremath{\eta}0=\ensuremath{\eta}(1keV), 对称且T(\ensuremath{\ell})=0, 求T(y).
\end{exercise}
\begin{solution}
(a)随机步长\ensuremath{\rho}与频率\ensuremath{\nu}给\ensuremath{\chi}\ensuremath{\perp}~\ensuremath{\nu}\ensuremath{\rho}\ensuremath{{}^{2}}, 按表
\begin{equation}
 \chi_{\perp e}\simeq1.580\times10^{-23}\frac n{B^2\sqrt{T_e}},\quad
 \chi_{\perp i}\simeq5.221\times10^{-22}\frac n{B^2\sqrt{T_i}}\quad(\mathrm m^2/\mathrm s).
\end{equation}
同T离子约33倍主导. (b)自由程v/\ensuremath{\nu}给\ensuremath{\chi}\ensuremath{\parallel}~v\ensuremath{{}^{2}}/\ensuremath{\nu}, 电子 $1.282\times10^{29}T_e^{5/2}/n$, 质子 $3.819\times10^{27}T_i^{5/2}/n$ m\ensuremath{{}^{2}}/s, 同T电子约34倍主导. 原表$v_t$按$\sqrt{T/m}$、$\rho$按$\sqrt{2T/m}/\Omega$, 有$\sqrt2$惯例差, 此处按题表分别估量级, 不替换系数而称精确.

(c) $(\alpha T^{-1/2}T')'+\eta_0J^2T^{-3/2}=0$. $u=\sqrt T$, 得 $2\alpha u''+\eta_0J^2u^{-3}=0$. 乘u\ensuremath{{}^{\prime}}积, $u^{\prime}(0)=0$、$u_0=\sqrt{T_0}$, 得 $u'^2=C(u^{-2}-u_0^{-2})$, C=\ensuremath{\eta}0J\ensuremath{{}^{2}}/(2\ensuremath{\alpha}). 再积分给 $T=T_0-Cy^2/T_0$, 边界得
\begin{equation}
 T_0=\ell|J|\sqrt{\frac{\eta_0}{2\alpha}},\qquad
 T(y)=T_0\left(1-\frac{y^2}{\ell^2}\right),\quad \eta_0=2.5\times10^{-8}\,\Omega\mathrm m.
\end{equation}
T为keV数值, \ensuremath{\alpha}已含能量单位转换. 到$T=0$边界热通量与电阻模型奇异, 故理想化抛物线不精确描述真实壁鞘层.
\end{solution}
\begin{exercise}\label{ex:06-final-3}
原题3. 高频波沿C-Mod中平面径向入射, 210GHz, 离子响应忽略. r用cm, $B_\varphi=5/(1+r/67)$T, $n_e=n_0[1-(r/20)^2]$. 所给nref\ensuremath{{}^{2}}图在r20为1、r12.5为0、r7.25为\ensuremath{\pm}\ensuremath{\infty}. (a)求n0; (b)r12.5极化(线、圆、椭圆, 平面和旋向); (c)用 $E''+(a/x+b)E=0$, x从共振向外测、E竖直分量, 求a/b; (d)给\ensuremath{\Gamma}=$1-\exp(-\pi\sqrt b\Delta)$, \ensuremath{\Delta}截止共振距离, 判断加热实用性.
\end{exercise}
\begin{solution}
(a)由X支外侧R截止 $\omega_{pe}^2=\omega(\omega-\Omega_{ce})$, 得 $n_0=3.9347\times10^{20}$m$^{-3}$. 若改从7.25cm的S0共振 $\omega_{pe}^2+\Omega_{ce}^2=\omega^2$反推, 得 $4.0200\times10^{20}$m$^{-3}$, 相差2.17\%. 原给两位置与剖面不能同时精确满足; 用(a)密度预测共振6.865cm. 故按原图作近似估量, 明示两约束冲突, 不捏造共同精确根.

(b)R0的介电本征矢量 $\mathbf E\propto\mathbf e_x+\ii\mathbf e_y$ (B沿局部z, e$^{-i\ensuremath{\omega}t}$), 是径向-竖直平面的圆偏振, 实场从+x向+y转, 与电子正B回旋同向. 截止k0时不再是纯横平面波, 径向E可非零. (c)沿原图几何取xcut=.0525m, xvac=.1275m, 模型 $k^2=a/x+b$在cut0、vac k0\ensuremath{{}^{2}}, 得
\begin{equation}
 b=k_0^2\frac{.1275}{.1275-.0525}=3.2931\times10^7\,\mathrm m^{-2},\quad
 a=-b(.0525)=-1.7289\times10^6\,\mathrm m^{-1}.
\end{equation}
这是拟合所给三位置的Budden局部模型, 不是整个密度剖面的精确介电函数. (d)$\pi\sqrt b\Delta\simeq946.48$, \ensuremath{\Gamma}近1, 实际几乎反射全部, 隔离层使此入射方案效率极差. \ensuremath{\Gamma}究竟振幅或题定义的系数依给定公式解释, 两者在该极限均近全反射. 本题不凭无碰撞冷奇点宣称得到实际加热率.

\begin{center}\begin{tikzpicture}\begin{axis}[width=.67\textwidth,height=4.6cm,xlabel={$r\ (\mathrm{cm})$},ylabel={$n_{\rm ref}^2$},xmin=0,xmax=20,ymin=-4,ymax=4,xtick={0,7.25,12.5,20},ytick={-2,0,1,2},unbounded coords=jump]
\addplot[thick,structurecolor,domain=0:7.1,samples=70]{1.7*(1-5.25/(x-7.25))};
\addplot[thick,structurecolor,domain=7.4:20,samples=120]{1.7*(1-5.25/(x-7.25))};
\addplot[dashed,domain=-4:4] coordinates{(7.25,-4)(7.25,4)};
\addplot[only marks]coordinates{(12.5,0)(20,1)};
\end{axis}\end{tikzpicture}\end{center}
按原题三位置与Budden拟合重绘的定性图: 左侧传播, 共振右侧至截止为倏逝, 截止之外可传播. 不将此曲线冒称给定密度剖面的精确冷介电结果.
\end{solution}
\begin{exercise}\label{ex:06-final-4}
原题4. 漂移电子 $f_e=n_0v_e/[\pi((v_x-V)^2+v_e^2)]$, 离子 $f_i=n_0v_i/[\pi(v_x^2+v_i^2)]$, Zi1, 波e$^{-i\ensuremath{\omega}t+ikx}$、k实\ensuremath{\omega}复. (a)代数色散; (b)题印低频准中性n1e+n1i=0重求; (c)解该式\ensuremath{\omega}(k); (d)漂移何时不稳定? 原件附Lorentz积分表.
\end{exercise}
\begin{solution}
(a) $1-\omega_{pe}^2/(\omega-kV+\ii|k|v_e)^2-\omega_{pi}^2/(\omega+\ii|k|v_i)^2=0$. (b)粒子数准中性应n1e=n1i; 源加号只在两个量定义为\textbf{带符号电荷密度}时成立, 与题的数密度定义不符. 因此采用消Poisson真空项1的长波式. (c)k>0、$\delta=\sqrt{m_e/m_i}$,
\begin{equation}
 \omega=\frac{k[\delta^2V+s\delta(v_i-v_e)]}{1+\delta^2}
 -\ii\frac{k[s\delta V+v_i+\delta^2v_e]}{1+\delta^2},\quad s=\pm1.
\end{equation}
负$k$的实系统对称为$\omega(-k)=-\omega(k)^*$, 或直接保留完整$|k|$式求解, 阻尼虚部不翻号. (d)存在\ensuremath{\gamma}>0需 $|V|>v_i/\delta+\delta v_e$, 冷离子阈\ensuremath{\delta}ve. 推导见章\ref{ch:instability}; 所有阻尼/增长根使用同一因果延拓, 不把下半平面普通实轴公式替换后产生虚假的相反符号.
\end{solution}
